\section{Reproducibility, Scope, and Operational Rule}
\label{sec:discussion}

\subsection{RQ9: Can the Decision Chain Be Replayed?}
The anonymous package contains 107 files and no identity leaks. A clean environment imports no main-repository code, rebuilds the paper and supplement, regenerates six figures and five data-linked tables, and passes 244 baseline checks, 1,205 post-seal checks, and 210/210 source-policy checks. A fresh Faiss environment reproduces registered native action rows with zero top-$k$, recall, NDC, or action-order mismatch. S9-3 adds 144,000 response cells, 112,000 timing measurements, and 16,000 top-$k$ checks with zero mismatch. S9-4 adds another 144,000 response cells and 20,500 checks, again with zero mismatch. Role manifests retain zero overlap among selection, certification, and evaluation.

The latest S9 records are committed with checksums and deterministic regeneration scripts. The procedural consistency audit follows a local scientific-writing checklist~\cite{kassis2026scientific}; all numeric claims remain bound to repository records rather than generated prose.

\subsection{What the Evidence Establishes}
Budget response changes across registered graph rebuilds, data refresh, implementations, and scale. ICBA detects this change and converts a frozen candidate into a target-qualified completed decision. TCP recovers endpoint-relative work for recurring profiled queries and adds resolved value beyond target-global calibration on SIFT. Target-global calibration is the stronger default on Arxiv. Source-derived slack can work, but S9-4 shows that a fixed action can match or dominate it.

The resulting maintenance rule has three steps. First, audit inherited policies against the target event and endpoint. Second, certify the least complex viable candidate and its fallback on disjoint target queries. Third, escalate from fixed to target-global to query-conditional policies only when held-out gain and lifecycle accounting justify the additional information.

\subsection{Boundaries}
The main conclusions are conditional on registered query populations, grids, SLAs, and build panels. The prospective S9 panel has eight builds per dataset and one CPU configuration; S9-4 reuses that build panel for paired attribution. TCP assumes recurring profiled queries. The complete cold-history ledger is expressed in NDC, while actual timing covers serving search rather than truth acquisition, control, energy, or storage. DARTH, Ada-ef, Deep1M, and Vamana retain different native estimands and support scope rather than a common ranking. These boundaries define the next replication targets without changing the fixed-target decisions reported here.

\section{Conclusion}
Search-budget portability is an index-conditioned deployment decision. ICBA makes that decision auditable by separating candidate construction, target certification, fallback, held-out evaluation, and cost. The evidence rejects automatic reuse, confirms safe recovery on fresh builds and queries, and shows why strong simple baselines are essential: fixed-256 can dominate source-tailored slack, while target calibration provides the largest robust gains and TCP is valuable when recurring query histories add measurable information. The practical outcome is a complexity-aware policy ladder backed by target evidence rather than a universal recovery rule.
